\documentclass[../main.tex]{subfiles}

\graphicspath{{\subfix{../images/}}}



\begin{document}


\subsection{Structure of a mains power supply}


\subsubsection{Conventional power supply}


% HERE PUT THE WHOLE INTPUT EMC TRANF RECT  FILTER


\subsubsection{Design example}
\subsubsection{Switching power supply}
\subsubsection{Design example}

\subsection{Power factor}



\subsubsection{Correct power factor displacement}



\subsection{Linear voltage regulators}


The purpose of a voltage regulator is to provide a stable output voltage to the circuit being powered, regardless of variations in the current drawn by the load and the input voltage.\\

Voltage regulators can be divided into two categories: linear (or dissipative) regulators and switching regulators.

Linear regulators always deliver an output voltage lower than the input voltage and are conceptually variable resistors, that is, resistors in which one of the two resistors changes its value to counteract variations in the load and/or the line.\\

To regulate the output voltage, one can conceptually envision using a divider composed of elements capable of changing their resistance depending on the voltage values present at the input and output, so as to keep the latter as constant as possible \autoref{linreg}.\\

The regulator must therefore be able to account for variations in both the load and the line. Since it makes no sense to change the values of both resistors simultaneously, there are two possible designs for the regulators:

\begin{enumerate}[label=\textbf{-}]
	\item Change the value of the divider element in parallel with the load, $R_P$ (parallel regulator);
	\item Apply the adjustment to the element in series with the $R_S$ load rather than to the one in parallel (series regulator); in this case, the parallel element can be eliminated, treating the load itself as the RP element of the divider.
\end{enumerate}




\begin{figure}[h!]
\begin{center}

\begin{tikzpicture}[transform shape]
	\node[shape=rectangle, draw, minimum width=3.986cm, minimum height=3.986cm] at (8, 6){};
	\draw (11, 7) to[generic, l={$R_L$}] (11, 5);
	\draw (8.625, 7) to[variable american resistor, l={$R_s$}] (5.625, 7);
	\draw (8.625, 7) to[variable american resistor, l={$R_p$}] (8.625, 5);
	\draw (11, 5) |- (5, 5);
	\draw (5.625, 7) -- (5, 7);
	\draw (8.625, 7) -- (11, 7);
	\node[ocirc] at (5, 7){};
	\node[ocirc] at (5, 5){};
\end{tikzpicture}

\label{linreg}
\caption{Generic circuit diagram of a voltage regulator}	
\end{center}	
\end{figure}



\subsubsection{Parallel regulator}

Given now a generic introduction to voltage regulators we can move on and dive deep inside some of the more specific types and configurations.\\

The simplest example of a parallel voltage regulator is constructed using a Zener diode as the parallel element \autoref{lindiode}.


\begin{figure}[h!]
\begin{center}

\begin{tikzpicture}[transform shape]
	\node[shape=rectangle, draw, minimum width=3.986cm, minimum height=3.986cm] at (8, 6){};
	\draw (11, 7) to[generic, l_={$R_L$}, v^<=$V_{L}$, voltage/distance from node=0.3, voltage/bump b=1.3, voltage/shift=2] (11, 5);
	\draw (11, 5) |- (5, 5);
	\draw (5.625, 7) -- (5, 7);
	\draw (8.625, 7) -- (11, 7);
	\node[ocirc] at (5, 7){};
	\node[ocirc] at (5, 5){};
	\draw (5.25, 7) to[american resistor, l={$R_S$}, i=$I_L$, current/distance=0.8] (9.125, 7);
	\draw (8.625, 5) to[empty ZZener diode, l={$D_Z$}] (8.625, 7);
	\draw (5, 5) to[open, v^=$V_{in}$, voltage/bump b=0, voltage/shift=-0.3] (5, 7);
\end{tikzpicture}
\label{lindiode}
\caption{Circuit diagram of a voltage regulator based on a Zener diode}	
\end{center}	
\end{figure}




If the Zener diode is properly biased, the voltage across its terminals remains approximately constant. 
The diode must be reverse-biased, and a current flowing through the diode must be at least greater than the minimum current $I_{Z_{min}}$. Under these conditions, the Zener characteristic is nearly linear, as shown in \autoref{zener}. 
The equivalent circuit of the Zener diode is, as a first approximation, an ideal voltage source with a voltage of $V_{Z0}$ in series with a resistance $R_Z$.

\begin{figure}[h!]
\begin{center}




\label{zener}
\caption{Reverse-biased characteristic of a Zener diode. The slope of the linear region defines $R_Z$.}	
\end{center}	
\end{figure} 

It is necessary to ensure, however, that the current does not exceed a maximum value $I_{Z_{max}}$, which would cause the Zener diode to dissipate excessive power.

Unfortunately, the series impedance of the Zener diode, $R_Z$, depends heavily on the current flowing through the diode. 

Above the typical knee point of the Zener’s characteristic curve, the impedance is inversely proportional to the current. Manufacturers generally define the nominal Zener voltage VZ as the voltage across the diode when it is traversed by a test current $I_{ZT}$, which is normally chosen to be around 30\% of the maximum current. The impedance value $Z_{ZT}$ is also provided for this current. 

For small variations in current around $I_{ZT}$, we can consider the impedance to be constant; however, if the current were significantly different, we would have to calculate the impedance value using the formula:

\begin{equation}
	Z_Z = Z_{ZT} \cdot \frac{I_{ZT}}{I_{ZT}}
\end{equation}


This formula provides an accurate value for currents starting at one-tenth of the $I_{ZT}$ current and not too close to the maximum current that the diode can handle—a range in which a series of parasitic impedances of the junction and the package become significant.

The voltage across the Zener diode deviates from $V_Z$ if the current through the diode is different from $I_{ZT}$. For small deviations, we can calculate $\Delta V_Z$ as:

\begin{equation}
	\Delta V_Z = \Delta I_Z Z_{ZT}
\end{equation}


If, on the other hand, we want to calculate the change for currents that are very different from $I_{ZT}$, then we must take into account the different differential impedance:

\begin{equation}
	\Delta V_Z = 2Z_{ZT} \cdot I_{ZT} \frac{I_Z - I_{ZT}}{I_Z + I_{ZT}}
\end{equation}

That said, normally, if it is not necessary to know the voltage across the diode with high precision, we proceed by using a simplified model in which the differential impedance is considered constant. 

The model is shown in \autoref{zenerss} . The value of $R_Z$ is taken from the component’s datasheet as $Z_{ZT}$, while the voltage $V_{Z0}$ is calculated as:

\begin{equation}
	V_{Z0} = V_Z - I_{ZT}Z_{ZT}
\end{equation}

that is, it is the value at which the linear approximation of the characteristic curve intersects the x-axis.


\begin{figure}[h!]
\begin{center}

\begin{tikzpicture}[transform shape]
	\draw (9, 5.375) to[empty ZZener diode, l={$D_Z$}] (9, 8);
	\draw (12, 8) to[american resistor, l={$R_Z$}] (12, 6.25);
	\draw (12, 6.25) to[battery1, l={$V_{Z0}$}] (12, 5.375);
	\draw (10.125, 6.625) -- (10.75, 6.625) -- (10.875, 6.625);
	\draw (10.125, 6.5) -| (10.875, 6.5);
	\draw (10.125, 6.375) -| (10.875, 6.375);
\end{tikzpicture}

\label{zenerss}
\caption{Equivalent circuit of a Zener diode}	
\end{center}	
\end{figure} 


The Zener voltage also depends on the junction temperature and on the aging of the device.


For special applications, there are two-pin active circuits capable of approximating the behavior of an ideal Zener diode (super-Zener) over a limited range of currents. 

Internally, they consist of a temperature-compensated reference generator that often uses a circuit known as a bandgap reference. 

We will not discuss them in detail; we will simply note that such circuits are typically used as voltage references in more complex circuits, whether they are voltage regulators or analog-to-digital conversion systems.









\subsubsection{Series Regulator}
\subsubsection{LDO Regulator}
\subsubsection{Design example of a power supply with a series regulator}







\end{document}